\section{Group or moment? The same flight before, during and after}\label{sec:prepost}

Do pilots fly together in the moments when the progress is straight anyway? If so, the same flight
would already be straight just before the group forms, and still just after it breaks.

\paragraph{Windows.} Every stretch of group time is an \emph{event}, with up to four windows of length
$W=\PPWindowMin$~min on the clock of the same flight: \emph{start} and \emph{end}, the first and last
$W$ of the group, when it lasts at least $W$; \emph{before} and \emph{after}, the last $W$ of the solo
stretch that precedes the group and the first $W$ of the one that follows it, when that solo stretch
lasts at least $W$, whatever the length of the group. A window of 50~min holds about 200 origins at
$\tau=1000$~s. For two windows the ensemble MSD is taken on the events that have both, and their
difference is read on $\Hloc$ and on the phase statistics. Entering is start minus before; leaving is
end minus after, group minus alone in both cases. The two readings predict different things:
\begin{itemize}
  \item if the straighter progress belongs to the group, it appears on entering and on leaving, and
        before and after are alike;
  \item if it belongs to the moment, before is already high, and the changes are those of any two
        consecutive windows.
\end{itemize}

\paragraph{Controls.} Two consecutive windows with no change of state, drawn in the solo time of the
group flights ($\A$) and in solo flights ($\Sl$), measure how much a flight drifts from one window to
the next, later minus earlier. Entering has the same time order as the controls; leaving has the
opposite one, so on leaving the drift enters with the opposite sign.

\begin{figure}
\centering
\includegraphics[width=\textwidth]{report_prepost.pdf}
\caption{Group minus alone on entering (start $-$ before) and on leaving (end $-$ after), each above the
drift it is read against: the controls later minus earlier for entering, earlier minus later for
leaving. Last row: after $-$ before. The curves over every lag are in Figure~\ref{fig:prepost}, the values in
Tables~\ref{tab:prepost} and~\ref{tab:prepost-phase}, the four windows in
Figure~\ref{fig:prepost-profile}.}
\label{fig:prepost-forest}
\end{figure}

\paragraph{Results} (Figure~\ref{fig:prepost-forest}). The first prediction holds at long lags:
$\Delta\Hloc(1000~\mathrm{s})$ is \PPEntryDHloc{} entering and \PPExitDHloc{} leaving, against drifts
near zero, and after minus before is \PPAfterBeforeDHloc: a step close to the within-flight estimate
of Section~\ref{sec:other}. At 100~s entering changes nothing beyond the drift, and on the fitted $H$
it exceeds the drift of $\A$ but not that of $\Sl$; on leaving the group window stays above both
drifts, at 100~s and on $H$: a small asymmetry, which we do not interpret. Glides are faster and flatter in the group both ways, by similar amounts once each drift
is taken with its sign. The climb statistics rise on entering beyond the drift and do not fall on
leaving beyond it.

\paragraph{Limits.} The windows fix the pilot, the glider, the day and the slow drift of the flight,
not what brings pilots together: a convergence line, a ridge or a common route may last as long as
the group and straighten the progress by itself, and two pilots must move alike to stay within
500~m for long. The events are long groups with solo flight beside them (\PPEntryEvents{} entries,
\PPExitEvents{} exits). With windows of \PPShortWindowMin~min the step at 1000~s is much smaller,
since such a window holds fewer than thirty origins at that lag; the test at that lag needs windows of
\PPWindowMin~min (Appendix~\ref{app:prepost}).

\paragraph{What we keep.} With windows of \PPWindowMin~min, on the same flight the straighter
progress over a quarter of an hour starts with the group and stops with it, and the glides change
with it both ways; the moment of joining does not explain them. The climbs change when the group
forms and not when it breaks, so for them the moment remains a possible cause.
